\section{割与覆盖参数：网络脆弱性不只看是否连通}
\label{sec:04-Discrete-Mathematics/03-Graph-Theory/11}

值班面板上写着``网络连通''，绿灯。第二天上午一台汇聚交换机例行维护下线，半个园区失联了四十分钟。事后复盘时那个绿灯还挂在墙上，它并没有说错——图确实连通，任意两点之间确实存在路径。它只是从来没有回答过被问的问题。

``连通''是一个布尔值。它把``删掉一个设备就散架''和``拆掉十条链路还能撑住''归成同一个答案。要区分这两种网络，需要的不是判断而是数字。

\subsection{离断开还有多远}

对连通的非完全图，顶点连通度

\[
\kappa(G)=\min\set{\abs{S}:G-S\text{ 不连通}}
\]

记录最少删掉多少个顶点会让图断开；边连通度

\[
\lambda(G)=\min\set{\abs{F}:G-F\text{ 不连通}}
\]

对边做同样的计量。若最小度为 \(\delta(G)\)，三者被一条链串起来：

\[
\kappa(G)\le\lambda(G)\le\delta(G).
\]

右边那个不等式几乎是显然的——挑一个度数最小的顶点，把它的关联边全删掉，它就孤立了。左边那个要深一些，但方向可以记住：破坏顶点比破坏边更省力，因为删一个顶点等于顺带删掉它的全部关联边。

\begin{center}
\begin{tikzpicture}[x=1cm,y=1cm,font=\small]
  \foreach \p in {(0,1.0),(0.9,1.9),(1.7,0.5),(0.5,-0.1)}{\fill \p circle (2.4pt);}
  \draw[black!70] (0,1.0) -- (0.9,1.9) -- (1.7,0.5) -- (0.5,-0.1) -- (0,1.0);
  \draw[black!70] (0,1.0) -- (1.7,0.5);
  \draw[black!70] (0.9,1.9) -- (0.5,-0.1);
  \foreach \p in {(4.9,1.4),(5.9,2.0),(6.5,0.5),(5.2,-0.2)}{\fill \p circle (2.4pt);}
  \draw[black!70] (4.9,1.4) -- (5.9,2.0) -- (6.5,0.5) -- (5.2,-0.2) -- (4.9,1.4);
  \draw[black!70] (4.9,1.4) -- (6.5,0.5);
  \draw[black!70] (5.9,2.0) -- (5.2,-0.2);
  \draw[very thick,red!70!black] (0.9,1.9) -- (4.9,1.4);
  \draw[very thick,red!70!black] (1.7,0.5) -- (5.2,-0.2);
  \draw[dashed,black!60,thick] (3.2,-1.1) -- (3.2,2.7);
  \node[black!60,font=\footnotesize] at (3.2,2.95) {割};
  \node[red!70!black,font=\footnotesize] at (1.85,2.3) {跨界边};
\end{tikzpicture}

\emph{两团内部各自很密，只有两条边跨过虚线。删掉这两条边，网络断成两半，所以 \(\lambda=2\)；与此同时，两团之间最多也只能找出两条互不共享边的路径。同一个数字有两个读法：最少要破坏多少，以及最多有几条独立的备用路。}
\end{center}

这两个读法总是给出同一个数，不是这张图的巧合。Menger 定理断言：两个不相邻顶点之间互不相交的路径的最大条数，恰好等于分开它们所需删除的顶点（或边）的最小个数。破坏方看到的最小代价，与建设方看到的最大冗余，是同一个量的两面。

图上的对偶还不止这一处。给边配上容量，最大流最小割定理说：从源点到汇点能推送的最大流量，等于所有分离源汇的割中容量最小的那一个。任何一个割的容量都是流量的上界——流从源到汇必须整个穿过这个割；而最小割恰好把这个上界压到可以达到。于是一份流方案和一个割，两者的数值一旦相等，就同时证明了流是最大的、割是最小的，谁也不用再搜索。这套``可行解给上界、对偶对象给下界，相等即最优''的手法，在\hyperref[ch:075]{``对偶：为最优性建立下界证书''}中会以一般形式重来一遍；上一节的旅行商近似里出现过它的组合版本。

最小割因此不只是脆弱性指标。既然它能把一张图切成内部紧密、跨界稀疏的两块，它也是一种聚类准则——只是朴素的最小割偏爱切下一个孤零零的顶点，实际使用的是归一化割之类带规模惩罚的版本，其连续松弛正是\hyperref[sec:13-Machine-Learning/05-Clustering-and-Data-Geometry/02]{``谱聚类把连通性变成特征向量''}的出发点。上面那张图里的虚线，换个语境读就是一条聚类边界。

模型的措辞必须写清楚删的是什么。基站故障对应删顶点，光缆被挖断对应删边，两者的连通度可以差很远。还有一个更常被忽略的前提：这些定义都假设故障相互独立地发生在不同元件上。若一批节点共用同一路市电、同一个机架或同一片洪泛区，它们会一起消失，此时 \(\kappa\) 给出的是一个乐观得离谱的数字。

\subsection{同一张图，四个不同的问题}

脆弱性只是网络设计要问的问题之一，换个问题就换一套参数，它们都长在同一张图上，彼此不能替换。

支配集 \(D\subseteq V\) 要求每个不在 \(D\) 里的顶点至少与 \(D\) 中某点相邻，问的是``充电桩、消防点或无线接入点放在哪几个节点，才能一跳覆盖全网''。点覆盖 \(C\subseteq V\) 要求每条边至少有一个端点落在 \(C\) 里，问的是``每一段关系至少由一端负责''。两个名字里都有``覆盖''，覆盖的对象却一个是顶点邻域、一个是边，混用会得到完全不同的答案。

点覆盖与独立集是同一件事的两种说法：\(C\) 是点覆盖当且仅当 \(V\setminus C\) 是独立集，于是最小点覆盖数与最大独立集数满足

\[
\tau(G)+\alpha(G)=\abs{V}.
\]

边覆盖与匹配则从相反方向使用边。边覆盖要每个顶点至少碰到一条选中的边，匹配要选中的边互不共享端点；前者怕有人被落下，后者怕资源被抢。若图没有孤立点，最小边覆盖数与最大匹配数满足

\[
\rho(G)+\nu(G)=\abs{V}.
\]

从一个最大匹配出发，给每个未匹配的顶点补一条关联边，就得到大小为 \(\abs{V}-\nu(G)\) 的边覆盖，反向论证给出最优性。\hyperref[sec:04-Discrete-Mathematics/03-Graph-Theory/04]{``二分图与匹配''}里的 König 定理在二分图上把这条线索接得更远：最大匹配数等于最小点覆盖数。

计算代价在这几个参数之间分布得极不均匀，这一点值得单独记住。最小割、最大流、最大匹配都有多项式算法，规模上万也能算；最小点覆盖和最小支配集在一般图上却是 NP 难的，只能靠近似、整数规划，或者利用二分图、树、有界树宽之类的额外结构。这又是一次上一节那个教训的重演：几个参数写在同一页纸上，措辞相近，落在计算世界的位置却隔着一整个阶级。

所以把需求叫作``网络优化''远远不够。动手之前先答三个问题：要覆盖的是顶点还是边，允许哪些元件共享风险，以及要抵抗的是哪一种失效。三个答案确定下来，参数才有唯一的选择；否则算出来的数字再精确，量的也不是你担心的那件事。

本单元从顶点与度数出发，一路走到匹配、遍历、可画性和割。它们共同的做法是把关于形状、位置、名字的问题，换算成只关于邻接关系的量。下一章换一个载体：把整数之间的整除关系当作结构来研究，会看到同样的思路——先找不变量，再看它能决定什么，\hyperref[ch:030]{``初等数论：整除结构怎样变成可计算算法''}从最大公约数开始。
